% Replacement for Section IV.C of the user-supplied CSCWD manuscript.
% Original conditions follow the user's Sun formula screenshots (2026-10-11).
% The user's preferred matching layout names the two original ratios.
% B, I, and S name original numerator sets; no decision condition is changed.

\subsection{Compliance Checking}
\label{sec:4.3}

Stage~3 combines the Process Record and Rule Record for rule--process matching and violation detection. We follow Sun et al.~\cite{sun2024design} and use their original notation in the formulas below. For rule \(r\), \(A_r\), \(P_r\), and \(U_r\) denote its actions, actors, and sequential constraints. For process model \(m\), \(A_m\), \(E_m\), and \(R_m\) denote its activities, events, and actors. The notation \(\mathit{bs\_obj}(A_m\cup E_m)\) denotes the business objects obtained from activity and event labels. The functions \(f_r\) and \(f_m\) and the relation \(F_m\) retain their meanings in Sun's definitions. The actor threshold \(\vartheta\) corresponds to \(\theta\) in our experimental setup.

\begin{figure*}[!b]
    \centering
    \includegraphics[width=0.98\textwidth,keepaspectratio]{fig2.png}
    \caption{Detailed workflow of Stage~2.}
    \label{fig:stage2}
\end{figure*}

\textbf{Rule--Process Matching.} The sets \(D_{r,m}\) and \(O_{r,m}\) contain the highest-similarity action and actor/business-object mapping pairs, respectively. To present Sun's matching formula compactly, \(m_A\) and \(m_O\) name its two original ratios. The score is displayed as \(\mathcal{F}\).

\begin{equation}
\begin{aligned}
m_A(r,m,\tau)
&=\frac{
  \displaystyle\sum_{(a_r,a_m)\in D_{r,m}}
  sim(a_r,a_m)>\tau
}{|D_{r,m}|},\\[2pt]
m_O(r,m,\tau)
&=\frac{
  \displaystyle\sum_{(o_r,o_m)\in O_{r,m}}
  sim(o_r,o_m)>\tau
}{|O_{r,m}|},\\[2pt]
\mathcal{F}(r,m,\tau)
&=\max\{m_A(r,m,\tau),m_O(r,m,\tau)\}.
\end{aligned}
\label{eq:matching}
\end{equation}

\textbf{Missing-Action Violation.} Let \(B_{r,m,\gamma}\) name the action-pair set in the numerator of Sun's original formula. Its condition and the denominator are retained.

\begin{equation}
B_{r,m,\gamma}=\{(a_r,a_m)\in D_{r,m}
  \mid sim(a_r,a_m)<\gamma\}.
\label{eq:missing_set}
\end{equation}

\begin{equation}
\mathit{action\_violation}(r,m,\gamma)
=\frac{|B_{r,m,\gamma}|}{|A_r|}.
\label{eq:missing}
\end{equation}

\textbf{Incorrect-Actor Violation.} The original rule-side and process-side candidate sets are reproduced first.

\begin{equation}
\begin{aligned}
R_{r,m,\gamma}
&=\bigl\{g\in P_r
  \mid\exists(a_r,a_m)\in D_{r,m},\\
&\qquad f_r(g)=a_r\land\\
&\qquad sim(a_r,a_m)>\gamma\bigr\}.
\end{aligned}
\label{eq:actor_rule_candidates}
\end{equation}

\begin{equation}
\begin{aligned}
C_{r,m,\gamma}
&=\bigl\{r\in\mathit{bs\_obj}(A_m\cup E_m)
  \cup R_m\mid\\
&\qquad\exists g\in R_{r,m,\gamma},\\
&\qquad sim(f_r(g),f_m(r))>\gamma\bigr\}.
\end{aligned}
\label{eq:actor_process_candidates}
\end{equation}

Let \(I_{r,m,\gamma,\vartheta}\) name the actor set in Sun's original numerator. The existential condition is retained without an extremum rewriting.

\begin{equation}
I_{r,m,\gamma,\vartheta}=\left\{
  \begin{gathered}
  r\in R_{r,m,\gamma}
  \mid\exists r'\in C_{r,m,\gamma},\\
  sim(r,r')<\vartheta
  \end{gathered}
  \right\}.
\label{eq:actor_mismatch_set}
\end{equation}

\begin{equation}
\mathit{actor\_violation}(r,m,\gamma,\vartheta)
=\frac{|I_{r,m,\gamma,\vartheta}|}{|R_{r,m,\gamma}|}.
\label{eq:actor}
\end{equation}

\textbf{Out-of-Order Violation.} The original evaluable sequential-constraint set is reproduced below.

\begin{equation}
\begin{aligned}
U_{r,m,\gamma}
&=\bigl\{(u_r,u'_r)\in U_r\mid\\
&\quad\exists\bigl\{(u_r,u_m)\in D_{r,m},\\
&\qquad (u'_r,u'_m)\in D_{r,m}\bigr\},\\
&\quad sim(u_r,u_m)>\gamma\land\\
&\quad sim(u'_r,u'_m)>\gamma\bigr\}.
\end{aligned}
\label{eq:order_evaluable}
\end{equation}

Let \(S_{r,m,\gamma}\) name the satisfied-constraint set in the numerator of Sun's original order-violation formula. Both endpoint mappings and both flow-relation conditions are retained.

\begin{equation}
S_{r,m,\gamma}=\left\{
  \begin{aligned}
  &(u_r,u'_r)\in U_{r,m,\gamma}\mid\\
  &\quad\exists\bigl\{(u_r,u_m)\in D_{r,m},\\
  &\qquad (u'_r,u'_m)\in D_{r,m}\bigr\},\\
  &\quad(u_m,u'_m)\in F_m\land\\
  &\quad(u'_m,u_m)\notin F_m
  \end{aligned}
  \right\}.
\label{eq:order_satisfied}
\end{equation}

\begin{equation}
\mathit{order\_violation}(r,m,\gamma)
=\frac{|U_{r,m,\gamma}|-|S_{r,m,\gamma}|}{|U_{r,m,\gamma}|}.
\label{eq:order}
\end{equation}

The ratios above apply when their denominators are nonzero. In our reconstruction, an empty matching set contributes a matching ratio of \(0\); a missing-action case with no rule action, an actor case with unavailable actor--action associations or candidates, and an ordering case with no evaluable relation produce no violation decision. An unmatched required action is treated as missing. These empty-input and unavailable-input policies are implementation conventions stated separately from Sun's formulas.
